A stargazer in Chiayi, Chinese Taipei ($23.5^\circ$N, $120.4^\circ$E, GMT+8)
saw two meteors streaking through the sky at 21:00 (Chinese Taipei time) on
25\textsuperscript{th} September 2020. One of the meteors appeared at horizon
exactly due west and streaked to a point at $15^\circ$ altitude directly above
the northern horizon. The second meteor originated at an altitude of
$23.5^\circ$ and an azimuth of $210^\circ$ and ended at an altitude of
$75^\circ$ and an azimuth of $255^\circ$.

\begin{description}
\item[(a)] (6 points) What is the Local Sidereal Time (LST) at the time of
  observation?
\item[(b)] (16 points) Find the alt-az coordinates of the apparent radiant of
  the two meteors.
\item[(c)] (6 points) Find the equatorial coordinates of the apparent radiant.
\item[(d)] (2 points) Which of the following constellations is closest to the
  radiant?\\
  Crux / Dorado / Pavo / Tucana / Triangulum Australes\\
  (choose one and write the same name in the answer box)
\end{description}

Notes:
\begin{itemize}
\item Azimuths are measured from the North ($0^\circ$) towards the East.
\item The Greenwich Sidereal Time (GST) at 00:00 UT on 1\textsuperscript{st}
  January 2020 is $6^{\mathrm{h}}\,40^{\mathrm{m}}\,30^{\mathrm{s}}$.
\end{itemize}
